\section{Proposed System and Key Features}

This section summarizes the system's inputs, key features, and outputs.

\subsection{System Concept}

Users can scan prescriptions, enter medication details manually, or submit requests through a chatbot or feature selection. OCR extracts prescription data, while the chatbot identifies multiple intents and explains system-provided results. After user verification and reliable medication identification, supported safety checks run automatically. Warnings are shown promptly; reminders follow the confirmed-schedule and safety-gate rules.

\subsection{System Inputs}

\subsubsection{Patient Input}

\begin{itemize}
    \item Printed or digital Vietnamese prescriptions, or manually entered medication information.

    \item User requests through feature selection or chatbot, including multiple intents.

    \item Relevant health information, reported symptoms, and confirmation of medication details and schedules.
\end{itemize}

\subsubsection{Healthcare Professional Input}

\begin{itemize}
    \item Professional assessments, schedule approvals, and clinical guidance.
\end{itemize}

\subsection{Key Features and Core Modules}

\subsubsection{Medication Intake \& Intent Identification}

This module extracts prescription information using OCR and accepts manual entries. A chatbot recognizes one or multiple intents and uses an existing language-model API to explain findings supplied by the safety and guidance modules, without generating independent safety conclusions.

\subsubsection{Medication Normalization}

This module matches user-verified medication details to reference records and identifies active ingredients. Unresolved identities are flagged, and checks involving those medications are marked incomplete rather than assumed safe.

\subsubsection{Medication Safety \& Risk Assessment}

This module automatically checks verified and reliably identified medications for drug-drug interactions and severity using DDInter 2.0. Results distinguish detected interactions, no known interaction found within available data, and incomplete or unassessed checks.

\subsubsection{Safety Guidance \& Professional Review}

This module immediately presents safety warnings, severity levels, and data limitations. It prepares high-risk or uncertain cases for professional review in the prototype and advises users to seek real clinical advice when needed.

\subsubsection{Medication Reminder \& Follow-up}

This module creates reminders from verified prescription schedules after user confirmation. Users may set notification times within prescribed instructions, while changes to dosing schedules require professional confirmation and user acceptance.

New reminders for medications involved in detected severe interactions are held until professional review and confirmation; unrelated reminders may proceed. The module also tracks taken or missed doses and records possible adverse drug reactions.

\subsection{System Outputs}

\begin{itemize}
    \item \textbf{Patient:} Verified medication information, safety findings, warnings, assessment limitations, medication reminders, adherence records, and professional guidance when available.

    \item \textbf{Healthcare Professional:} Structured patient summaries, medication information, detected risks, assessment limitations, and cases requiring professional review or schedule approval.
\end{itemize}
